\documentclass[a4paper, 12pt]{article}

% Essential Packages
\usepackage{amsmath}
\usepackage{amssymb}
\usepackage[top=25mm, bottom=25mm, left=25mm, right=25mm]{geometry}
\usepackage{pgfplots}
\usepackage[inline]{enumitem}
\usepackage{mathtools}
\usepackage{tikz}
\usepackage{fancyhdr}
\usepackage{setspace}
\usepackage{hyperref}

% Misc
\pgfplotsset{compat=1.18}
\usepgfplotslibrary{polar}
\usepgfplotslibrary{fillbetween}
\setlength{\parskip}{3pt}
\setstretch{1.15}
\renewcommand{\headrulewidth}{0pt}
\renewcommand{\footrulewidth}{0pt}
\fancyhf{}

\begin{document}
\pagestyle{empty}

\begin{center}
\textbf{2012-2013 Fall \\MAT123 Final\\23/01/2013}
\end{center}

\begin{enumerate}[leftmargin=*,label=\textbf{\arabic*.}]

\item Determine whether each given sequence with $n$th term converges or diverges. Evaluate the limit of each convergent sequence. Explain all your work, and write clearly.

\begin{enumerate}[leftmargin=*,label=\textbf{\alph*.}]
\item $\displaystyle a_n=(-1)^n\,n\sin\left(\frac1n\right)$
\item $\displaystyle a_n= e^{\cos\left(\textstyle\frac1n\right)}$
\end{enumerate}

\item Determine whether the following series converge or diverge. Give reasons for your answers.
\begin{enumerate}[leftmargin=*,label=\textbf{\alph*.}]
\item $\displaystyle\sum_{n=1}^\infty\frac n{\left(3+n^2\right)^{3/4}}$
\item $\displaystyle\sum_{n=2}^\infty\frac1{\sqrt n}\ln\left(\frac{n+1}{n-1}\right)$
\item $\displaystyle\sum_{n=1}^\infty\frac{(2n)!}{5^n\left(n!\right)^2}$
\item $\displaystyle\sum_{n=1}^\infty\frac1{n^2} e^{1/n}$
\end{enumerate}

\item Integrate the following functions and write each step in detail.
\begin{enumerate}[leftmargin=*,label=\textbf{\alph*.}]
\item $\displaystyle\int\frac{dx}{ e^x+1}$
\item $\displaystyle\int x\arcsin x\,dx$
\end{enumerate}

\item Find the length of the curve $\displaystyle y=\int_0^x\sqrt{\cos\left(4t\right)}\,dt$ for $0\leq x\leq \pi/8$.
\item For the function $\displaystyle f(x)=\frac x{x^2-4}$,
\begin{enumerate}[leftmargin=*,label=\textbf{\alph*.}]
\item Find all the asymptotes of $f$.
\item Find the intervals of increase or decrease.
\item Find the local maximum and minimum values, if any.
\item Find the intervals of concavity and the inflection points, if any.
\item Sketch the graph of $f(x)$.
\end{enumerate}
\end{enumerate}

\newpage
\pagestyle{fancy}
\renewcommand{\headrulewidth}{1pt}
\fancyhead[R]{\textit{2012-2013 Fall Final}}
\fancyhead[L]{\textit{Hacettepe MAT123 Exam Solutions Version 2}}

% Last update: 30/08/2025 00:29

\begin{enumerate}[leftmargin=*,label=\textbf{\arabic*.}]

\item \begin{enumerate}[leftmargin=*,label=\textbf{\alph*.}]
\item Evaluate the limit of the non-alternating part at infinity.
\[\lim_{n\to\infty}n\sin\left(\frac1n\right)\overset{n=\frac1u}{=}\lim_{u\to0}\frac{\sin u}u=1\]

However, for odd values of $n$, the limit at infinity becomes negative. On the other hand, for even values of $n$, the limit at infinity becomes positive. Therefore, the limit at infinity does not exist. So, the sequence diverges.
\item The exponential function $ e^x$ and the trigonometric function $\cos x$ are continuous everywhere.

\[\lim_{n\to\infty}a_n=\lim_{n\to\infty} e^{\displaystyle\cos\left(\frac1n\right)}= e^{\displaystyle\cos\left(\lim_{n\to\infty}\frac1n\right)}= e^{\cos0}= e^1= e\]
The sequence converges to $\boxed{ e}$.
\end{enumerate}

\item \begin{enumerate}[leftmargin=*,label=\textbf{\alph*.}]

\item Let $a_n=f(n)$. Define $\displaystyle f(x)=\frac x{\left(3+x^2\right)^{3/4}}$. The function is continuous for $x>1$ because the numerator and the denominator are continuous for $x>1$ and $\left(3+x^2\right)^{3/4}\neq0,\:\forall x\in\mathbb{R}$.
\[\left.\begin{array}{c}
x>0\\
\left(3+x^2\right)^{3/4}>0
\end{array}\right\}\:\text{for}\:x>1\implies\frac x{\left(3+x^2\right)^{3/4}}>0\]

For $x>1$, $x<x^{3/2}=\left(x^2\right)^{3/4}<\left(3+x^2\right)^{3/4}$. The denominator grows faster than the numerator. Therefore, the function is decreasing.

We may now apply the Integral Test. Handle the improper integral by taking the limit.
\begin{align*}
\int_1^\infty\frac{x}{(3+x^2)^{3/4}}\,dx
&= \lim_{R\to\infty}2(3+x^2)^{1/4}\bigg|_1^R\\[1em]
&= 2\lim_{R\to\infty}\left[(3+R^2)^{1/4}-(3+1^2)^{1/4}\right]= \infty
\end{align*}

Since the integral diverges, the series also $\boxed{\text{diverges}}$.

\item $\ln(1+x)<x$ for $x>-1$. Therefore,
\[\frac1{\sqrt n}\ln\left(\frac{n+1}{n-1}\right)=\frac1{\sqrt n}\ln\left(1+\frac2{n-1}\right)<\frac1{\sqrt n}\cdot\frac2{n-1}\]

Since $n\geq2$, we have the inequality $\displaystyle 2n-2\geq n\implies\frac2n\geq\frac1{n-1}$.
\[\frac1{\sqrt n}\cdot\frac2{n-1}\leq\frac1{\sqrt n}\cdot\frac4n=\frac4{n^{3/2}}\]

Now, let $\displaystyle a_n=\frac1{\sqrt n}\ln\left(\frac{n+1}{n-1}\right)$ and $\displaystyle b_n=\frac4{n^{3/2}}$. Apply the Direct Comparison Test.
\[0<a_n<b_n\implies0<\frac1{\sqrt n}\ln\left(\frac{n+1}{n-1}\right)<\frac4{n^{3/2}}\]

$b_n$ converges by the $p$-series test because $3/2>1$. Since $b_n$ converges, by the Direct Comparison Test, $a_n$ also $\boxed{\text{converges}}$.

\item Apply the Ratio Test. Let $\displaystyle a_n=\frac{(2n)!}{5^n\left(n!\right)^2}$.
\begin{align*}\lim_{n\to\infty}\left|\frac{a_{n+1}}{a_n}\right|&=\lim_{n\to\infty}\left|\frac{(2(n+1))!}{5^{n+1}\left((n+1)!\right)^2}\cdot\frac{5^n\left(n!\right)^2}{(2n)!}\right|\\[1em]&=\lim_{n\to\infty}\left|\frac{(2n+2)\cdot(2n+1)\cdot((2n)!)\cdot(n!)^2\cdot5^n}{5^n\cdot5\cdot(n+1)^2\cdot(n!)^2\cdot((2n)!)}\right|\\[1em]&=\lim_{n\to\infty}\left|\frac{(2n+2)(2n+1)}{5(n+1)^2}\right|=\lim_{n\to\infty}\left|\frac{4n^2+6n+2}{5n^2+10n+5}\right|\end{align*}

Now, take the corresponding function and evaluate the limit using L'Hôpital's rule.
\[\lim_{x\to\infty}\left|\frac{4x^2+6x+2}{5x^2+10x+5}\right|\overset{\text{L'H.}}{=}\lim_{x\to\infty}\left|\frac{8x+6}{10x+10}\right|\overset{\text{L'H.}}{=}\lim_{x\to\infty}\left|\frac{8}{10}\right|=\frac45<1\]

By the Ratio Test, the series converges absolutely. Since the series converges absolutely, the series $\boxed{\text{converges}}$.

\item Let $\displaystyle a_n=\frac1{n^2} e^{1/n}$. Define $\displaystyle f(x)=\frac1{x^2} e^{1/x}$. The function is continuous for $x>1$ because the expressions $\displaystyle\frac1{x^2}$ and $ e^{1/x}$ are continuous for $x>1$ and $x^2\neq0,\:\forall x>1$.
\[\left.\begin{array}{c}
x^2>0\\
 e^{1/x}>0
\end{array}\right\}\:\text{for}\:x>1\implies\frac{ e^{1/x}}{x^2}>0\]

For $x>1$, $ e^{1/x}$ tends to $1$ and $x^2$ is increasing. Therefore, the function is decreasing for $x>1$.

We may now apply the Integral Test. Handle the improper integral by taking the limit.
\[\int_1^\infty\frac{ e^{1/x}}{x^2}\,dx-=\lim_{R\to\infty}- e^{1/x}\bigg|_1^\infty=\lim_{R\to\infty}\left(- e^{1/R}+ e^{1}\right)= e-1\quad(\text{convergent})\]

Since the integral converges, by the Integral Test, the series also $\boxed{\text{converges}}$.
\end{enumerate}

\item \begin{enumerate}[leftmargin=*,label=\textbf{\alph*.}]
\item Add and subtract $ e^x$ in the numerator.
\[\int\frac{dx}{ e^x+1}=\int\frac{1+ e^x- e^x}{ e^x+1}\,dx=\int\frac{ e^x+1}{ e^x+1}\,dx-\int\frac{ e^x}{ e^x+1}\,dx\]

The result of the integral on the left is $x$. To calculate the integral on the right, use the $u$-substitution. Let $u= e^x+1$, then $du= e^x\,dx$.
\[x-\int\frac{du}{u}=x-\ln|u|+c=\boxed{x-\ln\left| e^x+1\right|+c=x-\ln\left( e^x+1\right)+c,\:c\in\mathbb{R}}\]

\item Solve the integral using the method of integration by parts.
\[\left.\begin{array}{c}
\displaystyle u=\arcsin x\implies du=\frac1{\sqrt{1-x^2}}\,dx\\[1em]
\displaystyle dv=x\,dx\implies v=\frac{x^2}2
\end{array}\right\}\rightarrow\int u\,dv=uv-\int v\,du\]

\begin{equation}\int x\arcsin x\,dx=\frac{x^2}2\arcsin x-\int\frac{x^2}{2\sqrt{1-x^2}}\,dx\label{eq:12_13_fall_final_1}\end{equation}

Now, we need to find the result of the integral on the right. Let $x=\sin u$ for $\displaystyle\frac\pi2<u<\frac\pi2$. Then $dx=\cos u\,du$. Also, recall the trigonometric identity $\sin^2x+\cos^2x=1$.
\begin{align*}\int\frac{x^2}{2\sqrt{1-x^2}}\,dx&=\int\frac{\sin^2u}{2\sqrt{1-\sin^2u}}\cdot\cos u\,du=\int\frac{\sin^2u\cos u}{2\sqrt{\cos^2u}}\,du\\[1em]&=\int\frac{\sin^2{u}\cos u}{2\left|\cos u\right|}\,du\quad\left[\cos u > 0\right]\\[1em]&=\frac12\int\sin^2u\,du=\frac12\int\left(1-\cos^2u\right)\,du=\frac12\int\left(\frac{1-\cos2u}2\right)\,du\\[1em]&=\frac u4-\frac{\sin2u}8+c=\frac u4-\frac{\sin u\cos u}4+c,\quad c\in\mathbb{R}\end{align*}

Recall the equation $x=\sin u$.
\[x=\sin u\implies\arcsin x=u\]
\[x=\sin u\implies x^2=\sin^2u=1-\cos^2u\implies\cos u=\sqrt{1-x^2}\]

Rewrite the result of the last integral.
\[\int\frac{x^2}{2\sqrt{1-x^2}}\,dx=\frac14\left(\arcsin x-x\sqrt{1-x^2}\right)\]

Rewrite \eqref{eq:12_13_fall_final_1}.
\[\int x\arcsin x\,dx=\boxed{\frac{x^2}2\arcsin x-\frac14\arcsin x+\frac14x\sqrt{1-x^2}+c,\quad c\in\mathbb{R}}\]

\end{enumerate}

\item The length of a curve defined by $y=f(x)$ whose derivative is continuous on the interval $a\leq x\leq b$ can be evaluated using the integral.
\[S=\int_a^b\sqrt{1+\left(\frac{dy}{dx}\right)^2}\,dx.\]

Find $\displaystyle\frac{dy}{dx}$.
\[\frac{dy}{dx}=\frac d{dx}\int_0^x\sqrt{\cos(4t)}\,dt\]

By the Fundamental Theorem of Calculus, $\displaystyle\frac{dy}{dx}$ can be rewritten as
\[\frac{dy}{dx}=\sqrt{\cos(4x)}\]

Set $\displaystyle a=0,\:b=\frac\pi8,\:\frac{dy}{dx}=\sqrt{\cos(4x)}$ and then find the length.

Recall $\cos(4x)=2\cos^2(2x)-1$.
\begin{align*}S&=\int_0^{\pi/8}\sqrt{1+\left(\sqrt{\cos(4x)}\right)^2}\,dx=\int_0^{\pi/8}\sqrt{1+\cos(4x)}\,dx\\[1em]&=\int_0^{\pi/8}\sqrt{2\cos^2(2x)}\,dx=\sqrt2\int_0^{\pi/8}\left|\cos(2x)\right|\,dx\quad\left[\cos(2x)>0\right]\\[1em]&=\sqrt2\int_0^{\pi/8}\cos(2x)\,dx=\frac{\sqrt2}2\sin(2x)\bigg|_0^{\pi/8}=\frac{\sqrt2}2\left(\sin\frac\pi4-\sin0\right)=\boxed{\frac12}\end{align*}

\item \begin{enumerate}[leftmargin=*,label=\textbf{\alph*.}]

\item Find the horizontal asymptotes.
\[\lim_{x\to\pm\infty}\frac{1}{x^2-4}=0\]

Find the vertical asymptotes. The expression is undefined for $x=\pm2$.
\[\lim_{x\to2^+}\frac{1}{x^2-4}=\lim_{x\to-2^+}\frac{1}{x^2-4}=\infty,\qquad \lim_{x\to2^-}\frac{1}{x^2-4}=\lim_{x\to-2^-}\frac{1}{x^2-4}=-\infty\]
\[\boxed{\text{The horizontal asymptote is } y= 0. \text{ The vertical asymptotes are } x=\pm2.}\]

\item Compute the first derivative and set it to $0$ to find the critical points. Apply the product rule appropriately.
\[f'(x)=\frac{1\cdot\left(x^2-4\right)-x\cdot(2x)}{\left(x^2-4\right)^2}=-\frac{x^2+4}{\left(x^2-4\right)^2}\]

$f$ is increasing where $f'(x)>0$ and decreasing where $f'(x)<0$. Therefore,
\[\boxed{\text{f is decreasing everywhere except at the undefined points.}}\]

\item $\boxed{\text{No local maximum or minimum values exist.}}$

\item Compute the second derivative.
\[f''(x)=-\frac{2x\cdot\left(x^2-4\right)^2-\left(x^2+4\right)\cdot2\cdot(x^2-4)\cdot(2x)}{\left(x^2-4\right)^4}=\frac{2x^3+24x}{\left(x^2-4\right)^3}\]
\[
\boxed{
\begin{array}{c}
\displaystyle \text{An inflection point occurs at }x=0.\\
\displaystyle f\text{ is concave up for }-2<x<0\:\cup\:x>2.\\
\displaystyle f\text{ is concave down for }x<-2\:\cup\:0<x<2.
\end{array}}
\]

\item The graph is below.
\begin{center}
\begin{tikzpicture}
\begin{axis}[
    domain=-6:6,
    ymin=-4, ymax=4,
    samples=500,
    axis lines=middle,
    xlabel={$x$},
    ylabel={$y$},
    restrict y to domain=-4:4,
    enlargelimits=true,
    grid=none,
    xtick={-6,-4,-2,0,2,4,6},
    ytick={-6,-4,-2,0,2,4,6},
    unbounded coords=jump,
    clip=true,
    scale=2
]
\addplot[blue, thick] {x/(x^2 - 4)};
\draw[dashed, red] (axis cs: -6,0.02) -- (axis cs: 6,0.02);
\draw[dashed, red] (axis cs: 2,-4) -- (axis cs: 2,4);
\draw[dashed, red] (axis cs: -2,-4) -- (axis cs: -2,4);

\end{axis}
\end{tikzpicture}
\end{center}
\end{enumerate}
\end{enumerate}
\end{document}